\section{Glossar}

Das Glossar definiert die Ubiquitous Language von EduPlanner.

\subsection{Fachliche Begriffe}

\begin{longtable}{L{3cm} L{2cm} p{\dimexpr\linewidth-5cm-6\tabcolsep\relax}}
  \rowcolor{tableheader}
  \color{white}\textbf{Begriff} & \color{white}\textbf{Kontext} & \color{white}\textbf{Definition} \\
  \endfirsthead
  \rowcolor{tableheader}
  \color{white}\textbf{Begriff} & \color{white}\textbf{Kontext} & \color{white}\textbf{Definition} \\
  \endhead
 Course & [catalog] & Weiterbildungsangebot mit Titel, Beschreibung und Lernzielen. Beschreibt \emph{was} gelehrt wird. Bsp.: ``Python für Einsteiger''. \\
  \rowcolor{tablegrey}
  CourseStatus & [catalog] & Zustand eines Kurses: DRAFT $\to$ PUBLISHED $\to$ ARCHIVED. Nur diese Übergänge erlaubt (State Machine). \\
  TrainerProfile & [catalog] & Stammdaten eines Trainers inkl.\ Fachgebiete. \\
  \rowcolor{tablegrey}
  CourseExecution & [scheduling] & Konkrete, zeitlich geplante Instanz eines Kurses. Bsp.: ``Python Apr~2025, Bern, max.\ 12~TN''. \\
  ExecutionStatus & [scheduling] & PLANNED $\to$ PUBLISHED $\to$ CANCELLED. Nur PUBLISHED offen für Einschreibungen. \\
  \rowcolor{tablegrey}
  Capacity & [scheduling, reg.] & Maximale Teilnehmeranzahl. Kapazität = 0 $\to$ Warteliste. \\
  Registration & [registration] & Verbindliche Anmeldung eines Teilnehmers zu einer Execution. Bsp.: Max Muster meldet sich an $\to$ \texttt{status=REGISTERED}. \\
  \rowcolor{tablegrey}
  WaitlistEntry & [registration] & Wartelisteneintrag mit Position (1 = erster) und Zeitstempel. \\
  Participant & [identity] & Natürliche Person, die sich einschreiben kann. Systemrolle: PARTICIPANT. \\
  \rowcolor{tablegrey}
  Trainer & [identity] & Kursleiter. Systemrolle: TRAINER. \\
  Staff & [identity] & Mitarbeitende des Kursanbieters. Systemrolle: STAFF. \\
  \rowcolor{tablegrey}
  Admin & [identity] & Nutzer mit erweiterten Systemrechten. Systemrolle: ADMIN. \\
  Domain Event & -- & Unveränderliche Zustandsänderung als Fakt (Vergangenheitsform). Bsp.: \texttt{ParticipantRegistered} (nicht \texttt{RegisterParticipant}). \\
  \rowcolor{tablegrey}
  Bounded Context & -- & Klar abgegrenzter Bereich mit eigener Fachsprache. Bsp.: ``Trainer'' in [catalog] hat andere Bedeutung als in [scheduling]. \\
  Aggregate & -- & Gruppe von Domain-Objekten als konsistente Einheit. Aggregate Root ist einziger Einstiegspunkt. Bsp.: Registration-Aggregate verhindert Überbuchung. \\
  \rowcolor{tablegrey}
  Ubiquitous Language & -- & Gemeinsame fachliche Sprache innerhalb eines BC. Code-Bezeichner verwenden exakt dieselben Begriffe wie die Fachsprache. \\
\end{longtable}

\subsection{Domain Events -- Vollständige Übersicht}

\begin{tabularx}{\textwidth}{L{5cm} L{2.5cm} X}
  \rowcolor{tableheader}
  \color{white}\textbf{Domain Event} & \color{white}\textbf{Kontext} & \color{white}\textbf{Auslöser} \\
  CourseCreated & catalog & Staff legt Kurs an \\
  \rowcolor{tablegrey}
  CourseTitleUpdated & catalog & Staff aktualisiert Kurstitel \\
  TrainerAssignedToCourse & catalog & Staff ordnet Trainer zu \\
  \rowcolor{tablegrey}
  CoursePublished & catalog & Admin veröffentlicht \\
  CourseArchived & catalog & Admin archiviert \\
  \rowcolor{tablegrey}
  CourseExecutionScheduled & scheduling & Staff plant Durchführung \\
  TrainerAssignedToExecution & scheduling & Staff weist Trainer zu \\
  \rowcolor{tablegrey}
  CourseExecutionPublished & scheduling & Admin gibt frei \\
  CourseExecutionCancelled & scheduling & Admin / System storniert \\
  \rowcolor{tablegrey}
  ParticipantRegistered & registration & Teilnehmer schreibt sich ein \\
  ParticipantAddedToWaitlist & registration & Kurs voll \\
  \rowcolor{tablegrey}
  RegistrationCancelled & registration & Teilnehmer / Staff storniert \\
  WaitlistParticipantPromoted & registration & Nachrücker von Warteliste \\
  \rowcolor{tablegrey}
  EmailNotificationSent & notification & System hat E-Mail versendet \\
  UserCreated & identity & Admin / IdP legt Benutzer an \\
  \rowcolor{tablegrey}
  UserRoleAssigned & identity & Admin weist Rolle zu \\
\end{tabularx}
